% \begin{definition}[Covering]
%     A covering map in $\Delta^n$ is a map $([a_1],\ldots, [a_n]) \to ([b_1], \ldots, [b_n])$ such that each component $[a_i] \to [b_i]$ preserves the minimum and maximum elements.
% \end{definition}\pablo{is that a standard definition? what's the standard name for that?}

% \begin{lemma}
%     Let $X$ be an $n$-fold simplicial space and $m_1, \ldots, m_n \geq 0$.
%     Let $x \inset \mathbb{L}_S(X)_{m_1, \ldots, m_n}$.
%     There exists a covering $([m_1],\ldots, [m_n]) \to ([c_1], \ldots, [c_n])$ such that $x$ lies in the image of
%     \[
%         X_{c_1, \ldots, c_n} \to \mathbb{L}_S(X)_{c_1, \ldots, c_n} \to \mathbb{L}_S(X)_{m_1, \ldots, m_n}
%     \]
% \end{lemma}
% \begin{proof}
%     Note that by \cite[Thm. D]{BS24}, the Segalification is obtained by successfully ``segalify'' the $k^\textnormal{th}$-coordinate, starting from the bottom up: $\mathbb{L}_S \simeq \mathbb{L}^n_S \circ \cdots \circ \mathbb{L}^1_S$.
%     Obvserve \cite[Obv. $1.16$]{BS24} that $\mathbb{L}^i_S(X)_{m_1, \ldots, m_{i-1}, 0, m_{i+1}, \ldots, m_n} \simeq X_{m_1, \ldots, m_{j-1}, 0, m_{j+1}, \ldots, m_n}$ for all $1 \leq i \leq n$.
%     We are trying to compute $\mathbb{L}^n_S \circ \cdots \circ \mathbb{L}^1_S (X)_{1, \ldots, 1, 0, \ldots, 0}$ with $k$ $1$'s, so we can ignore all $\mathbb{L}^i_S$ with $i > k$.
% \end{proof}

% \begin{lemma}
%     Let $1 \leq k < n$.
%     Then $\Psh(\Delta^n)^{k-\adj} \mono \Psh(\Delta^n)$ is stable under limits and colimits.
% \end{lemma}

% \begin{lemma}
%     Let $1 \leq k < n$.
%     Let $X$ be an $n$-fold simplicial space that has right (resp. left) adjoints for all $k$-cells.
%     Then its Segalification $\mathbb{L}_S(X)$ has right (resp. left) adjoints for all $k$-cells.
% \end{lemma}
% \begin{proof}
%     We proceed by induction on $k$.
%     \begin{enumerate}
%         \item \textbf{Base case:} Suppose $k = 1$.
%         We have $\mathbb{L}_S(X)_{1,0, \ldots, 0} \simeq (\mathbb{L}_S^n \circ \mathbb{L}_S^1 (X))_{1,0, \ldots, 0} \simeq \mathbb{L}_S^1 (X)_{1,0, \ldots, 0}$ by \cite[Thm. D]{BS24} and \cite[Obv. $1.16$]{BS24}.
%         Let $x \inset \mathbb{L}_S^1 (X)_{1,0, \ldots, 0}$.
%         Using the explicit necklace formula \cite[Thm. A]{BS24}, we can find $n_1, \ldots, n_r$ such that $x$ lies in the image of 
%         \[
%             X_{n_1,0, \ldots, 0} \times_{X_{0, \ldots, 0}}  X_{n_2,0, \ldots, 0} \times_{X_{0, \ldots, 0}} \cdots \times_{X_{0, \ldots, 0}}  X_{n_r,0, \ldots, 0} \longto \mathbb{L}_S^1 (X)_{1,0, \ldots, 0}
%         \]
%         Every element in $X_{n,0, \ldots, 0}$ has an adjoint by assumption, and adjoints are stable under composition \pablo{prove this for non-segal spaces}, so $x$ has an adjoint.

%         \item \textbf{Inductive step:} Suppose that $2 \leq k$ and that the result true for $k-1$.
%         Let $m_1 \geq 0$.
%         The necklace formula gives $\mathbb{L}_S^1 (X)_{m_1, \bullet, \ldots, \bullet}$ in terms of the $X_{k, \bullet, \ldots, \bullet}$ using some limits and colimits.
%         By assumption, each $X_{k, \bullet, \ldots, \bullet} \inset \Psh(\Delta^{(n-1)})^{(k-1)-\adj}$ and hence $\mathbb{L}_S^1 (X)_{m_1, \bullet, \ldots, \bullet} \inset \inset \Psh(\Delta^{(n-1)})^{(k-1)-\adj}$.
%         This exactly means that $\mathbb{L}_S^1 (X)$ has adjoints for all $k$-cells.
%         By the 
%     \end{enumerate}
% \end{proof}

% \begin{lemma}
%     As above, let $k + d \leq n$ and $\sigma = \sigma_{k+d-1} \ldots \sigma_{k} \inset \{\epsilon, \eta\}^d$ be a $d$-duality data.
%     \begin{enumerate}
%         \item If $\sigma_{k+d-1} = \epsilon$, then the ``domain'' weak constructible bundle $\left(\mathbb{R}^{n}_{\overline{k},\overline{\sigma}} \right)_{x_{k+d} < 0} \to \left(\mathbb{R}^{n}_{\overline{k+d-1},\overline{\sigma_{k+d-1}}}\right)_{x_{k+d} < 0} \to \left(\mathbb{R}^{n}_{\overline{k+d}}\right)_{x_{k+d} < 0}$ is isomorphic to a composition
%         \[
%             X \to  \mathbb{R}^{n}_{\overline{k+d} \circ \overline{k+d}} \to \mathbb{R}^{n}_{\overline{k+d}}
%         \]
%         where the first w.c. bundle, restricted to each of $x_{k+d-1} < 1$ and  $x_{k+d-1} > -1$, is isomorphic to $\mathbb{R}^{n}_{\overline{k},\overline{\sigma_{k+d-2}\ldots\sigma_k}} \to \mathbb{R}^{n}_{\overline{k+d-1}}$.
%         \item Similarly for $\sigma_{k+d-1} = \eta$ with the ``codomain'' $x_{k+d} > 0$.
%         \item If $\sigma_{k+d-1} = \epsilon$, then the ``domain'' weak constructible bundle $\left(\mathbb{R}^{n}_{\overline{k},\overline{\sigma}} \right)_{x_{k+d} < 0} \to \left(\mathbb{R}^{n}_{\overline{k+d}}\right)_{x_{k+d} < 0}$ is isomorphic to
%         \[
%             \mathbb{R}^{n}_{\overline{k},\overline{\sigma_{k+d-2}\ldots\sigma_k}} \to \mathbb{R}^{n}_{\overline{k+d-1}}
%         \]
%     \end{enumerate}
% \end{lemma}
% \begin{proof}
%     Note that, because the $x_{k+d}$-coordinate never changes, we have a commutative diagram where the downward maps are fiber bundles:
%     % https://q.uiver.app/#q=WzAsNSxbMCwwLCJcXGxlZnQoXFxtYXRoYmJ7Un1ee259X3tcXG92ZXJsaW5le2t9LFxcb3ZlcmxpbmV7XFxzaWdtYX19XFxyaWdodClfe3hfe2srZH0gPCAwfSJdLFsyLDAsIlxcbGVmdChcXG1hdGhiYntSfV57bn1fe1xcb3ZlcmxpbmV7aytkLTF9LFxcb3ZlcmxpbmV7XFxzaWdtYV97aytkLTF9fX1cXHJpZ2h0KV97eF97aytkfSA8IDB9Il0sWzMsMCwiXFxsZWZ0KFxcbWF0aGJie1J9XntufV97XFxvdmVybGluZXtrK2R9fVxccmlnaHQpX3t4X3trK2R9IDwgMH0iXSxbMSwwLCJcXGNkb3RzIl0sWzMsMSwiXFxtYXRoYmJ7Un1fezwwfSJdLFsxLDJdLFswLDNdLFszLDFdLFsyLDQsInhfe2srZH0iXSxbMSw0LCJ4X3trK2R9IiwxXSxbMCw0LCJ4X3trK2R9IiwyLHsiY3VydmUiOjJ9XSxbOSwxMCwiXFxjZG90cyIsMyx7InNob3J0ZW4iOnsic291cmNlIjoyMCwidGFyZ2V0IjoyMH0sInN0eWxlIjp7ImJvZHkiOnsibmFtZSI6Im5vbmUifSwiaGVhZCI6eyJuYW1lIjoibm9uZSJ9fX1dXQ==
% \[\begin{tikzcd}
% 	{\left(\mathbb{R}^{n}_{\overline{k},\overline{\sigma}}\right)_{x_{k+d} < 0}} & \cdots & {\left(\mathbb{R}^{n}_{\overline{k+d-1},\overline{\sigma_{k+d-1}}}\right)_{x_{k+d} < 0}} & {\left(\mathbb{R}^{n}_{\overline{k+d}}\right)_{x_{k+d} < 0}} \\
% 	&&& {\mathbb{R}_{<0}}
% 	\arrow[from=1-1, to=1-2]
% 	\arrow[""{name=0, anchor=center, inner sep=0}, "{x_{k+d}}"', curve={height=12pt}, from=1-1, to=2-4]
% 	\arrow[from=1-2, to=1-3]
% 	\arrow[from=1-3, to=1-4]
% 	\arrow[""{name=1, anchor=center, inner sep=0}, "{x_{k+d}}"{description}, from=1-3, to=2-4]
% 	\arrow["{x_{k+d}}", from=1-4, to=2-4]
% 	\arrow["\cdots"{marking, allow upside down}, draw=none, from=1, to=0]
% \end{tikzcd}\]
%     By \cite[Lem. $6.1.5$]{AFR18a}\pablo{almost it}, we get an identification of the horizontal pinch maps with a product of $\mathbb{R}$ (in $k+d$-position) and
%     % https://q.uiver.app/#q=WzAsNCxbMCwwLCJcXGxlZnQoXFxtYXRoYmJ7Un1ee259X3tcXG92ZXJsaW5le2t9LFxcb3ZlcmxpbmV7XFxzaWdtYX19XFxyaWdodClfe3hfe2srZH0gPSAtMX0iXSxbMiwwLCJcXGxlZnQoXFxtYXRoYmJ7Un1ee259X3tcXG92ZXJsaW5le2srZC0xfSxcXG92ZXJsaW5le1xcc2lnbWFfe2srZC0xfX19XFxyaWdodClfe3hfe2srZH0gPSAtMX0iXSxbMywwLCJcXGxlZnQoXFxtYXRoYmJ7Un1ee259X3tcXG92ZXJsaW5le2srZH19XFxyaWdodClfe3hfe2srZH0gPSAtMX0iXSxbMSwwLCJcXGNkb3RzIl0sWzEsMl0sWzAsM10sWzMsMV1d
% \[\begin{tikzcd}
% 	{\left(\mathbb{R}^{n}_{\overline{k},\overline{\sigma}}\right)_{x_{k+d} = -1}} & \cdots & {\left(\mathbb{R}^{n}_{\overline{k+d-1},\overline{\sigma_{k+d-1}}}\right)_{x_{k+d} = -1}} & {\left(\mathbb{R}^{n}_{\overline{k+d}}\right)_{x_{k+d} = -1}}
% 	\arrow[from=1-1, to=1-2]
% 	\arrow[from=1-2, to=1-3]
% 	\arrow[from=1-3, to=1-4]
% \end{tikzcd}\]
%     Note that the last map, once multiplied with $\mathbb{R}$ as described, is precisely the composition pinch map $\mathbb{R}^{n}_{\overline{k+d} \circ \overline{k+d}} \to \mathbb{R}^{n}_{\overline{k+d}}$.
%     Now it suffices to show that $\left(\mathbb{R}^{n}_{\overline{k},\overline{\sigma}}\right)_{x_{k+d} = -1, x_{k+d-1} > -1} \to \left(\mathbb{R}^{n}_{\overline{k+d-1},\overline{\sigma_{k+d-1}}}\right)_{x_{k+d} = -1, x_{k+d-1} > -1}$ is isomorphic to $\mathbb{R}^{n}_{\overline{k},\overline{\sigma_{k+d-2}\ldots\sigma_k}} \to \mathbb{R}^{n}_{\overline{k+d-1}}$.
%     \pablo{Write explicit rescaling}
% \end{proof}

% Finally, we identify the full $n$-duality datum as actions from higher spherical versions of Hoschild Homology.
% Note that we have a canonical closed inclusion of strata $\mathbb{R}_{\geq 0} \mono \mathbb{R}^n_{\overline{n}}$ and $\mathbb{R}_{\leq 0} \mono \mathbb{R}^n_{\overline{n}}$ on the last coordinate.
% When $M$ is a $\mathbb{R}^n_{\overline{n}}$-algebra in $\V$, we will often talk about the 

% \begin{lemma}[Action by Spherical Hochschild Homology]\label{lem:hosschild-action}
%     Let $k + d \leq n$ and $\sigma = \sigma_{k+d-1} \ldots \sigma_{k} \inset \{\epsilon, \eta\}^d$ be a $d$-duality data.
%     Let $a = |\{ k \leq i < k+d \mid \sigma_i = \epsilon\}|$ and $b = d - a$.
%     There is a commutative diagram
    
% \end{lemma}
% \begin{proof}
%     % https://q.uiver.app/#q=WzAsOCxbMiwwLCJcXGxlZnQoXFxtYXRoYmJ7Un1ebl97XFxvdmVybGluZXtrfSwgXFxvdmVybGluZXtcXHNpZ21hfX1cXHJpZ2h0KV97eF97aytkfSBcXGdlcSAwfSJdLFszLDAsIlxcbGVmdChcXG1hdGhiYntSfV5uX3tcXG92ZXJsaW5le2srZH19XFxyaWdodClfe3hfe2srZH0gXFxnZXEgMH0iXSxbMywxLCJDKFxcYnVsbGV0KSJdLFsyLDEsIlgiXSxbMSwwLCJcXGxlZnQoXFx7MFxcfSBcXGluc2V0IFxcbWF0aGJie1J9Xm5fe1xcb3ZlcmxpbmV7a30sIFxcc2lnbWF9XFxyaWdodClfe3hfe2srZH0gXFxnZXEgMH0iXSxbMSwxLCJcXHswXFx9IFxcaW5zZXQgXFxtYXRoYmJ7Un1eYSBcXHRpbWVzIFxcbGVmdChcXG1hdGhiYntSfV5uX3tcXG92ZXJsaW5le2t9LCAxXFxsZG90czF9XFxyaWdodClfe3hfe2srZC1hfSBcXGdlcSAwfSJdLFswLDEsIlxcezBcXH0gXFxpbnNldCBcXG1hdGhiYntSfV5hIFxcdGltZXMgXFxtYXRoYmJ7Un1fe1xcZ2VxIDB9Il0sWzAsMCwiXFx7MFxcfSBcXGluc2V0IFxcbWF0aGJie1J9XmEgXFx0aW1lcyBcXG1hdGhiYntSfV97XFxnZXEgMH0iXSxbMSwyLCIiLDAseyJzdHlsZSI6eyJoZWFkIjp7Im5hbWUiOiJlcGkifX19XSxbMCwzLCIiLDAseyJzdHlsZSI6eyJoZWFkIjp7Im5hbWUiOiJlcGkifX19XSxbMywyLCIiLDEseyJzdHlsZSI6eyJoZWFkIjp7Im5hbWUiOiJlcGkifX19XSxbMCwxLCIiLDEseyJzdHlsZSI6eyJoZWFkIjp7Im5hbWUiOiJlcGkifX19XSxbMCw0LCIiLDAseyJzdHlsZSI6eyJib2R5Ijp7Im5hbWUiOiJkb3R0ZWQifX19XSxbNCw1LCIiLDAseyJzdHlsZSI6eyJoZWFkIjp7Im5hbWUiOiJlcGkifX19XSxbMyw1LCIiLDAseyJzdHlsZSI6eyJib2R5Ijp7Im5hbWUiOiJkb3R0ZWQifX19XSxbNSw2LCIiLDAseyJzdHlsZSI6eyJoZWFkIjp7Im5hbWUiOiJlcGkifX19XSxbNyw2LCIiLDIseyJsZXZlbCI6Miwic3R5bGUiOnsiaGVhZCI6eyJuYW1lIjoibm9uZSJ9fX1dLFs0LDcsIiIsMix7InN0eWxlIjp7ImhlYWQiOnsibmFtZSI6ImVwaSJ9fX1dXQ==
% \[\begin{tikzcd}
% 	{\{0\} \inset \mathbb{R}^a \times \mathbb{R}_{\geq 0}} & {\left(\{0\} \inset \mathbb{R}^n_{\overline{k}, \sigma}\right)_{x_{k+d} \geq 0}} & {\left(\mathbb{R}^n_{\overline{k}, \overline{\sigma}}\right)_{x_{k+d} \geq 0}} & {\left(\mathbb{R}^n_{\overline{k+d}}\right)_{x_{k+d} \geq 0}} \\
% 	{\{0\} \inset \mathbb{R}^a \times \mathbb{R}_{\geq 0}} & {\{0\} \inset \mathbb{R}^a \times \left(\mathbb{R}^n_{\overline{k}, 1\ldots1}\right)_{x_{k+d-a} \geq 0}} & X & {C(\bullet)}
% 	\arrow[equals, from=1-1, to=2-1]
% 	\arrow[two heads, from=1-2, to=1-1]
% 	\arrow[two heads, from=1-2, to=2-2]
% 	\arrow[dotted, from=1-3, to=1-2]
% 	\arrow[two heads, from=1-3, to=1-4]
% 	\arrow[two heads, from=1-3, to=2-3]
% 	\arrow[two heads, from=1-4, to=2-4]
% 	\arrow[two heads, from=2-2, to=2-1]
% 	\arrow[dotted, from=2-3, to=2-2]
% 	\arrow[two heads, from=2-3, to=2-4]
% \end{tikzcd}\]
% \end{proof}

% \subsection{Enriched Absolute Colimits}

% Let $\V \in \EAlg{1}(\Cat{1}^{\Colim})$.
% In \cite{Hau16}, Haugseng defines a \textit{beta}-version of the $1$-Morita category as a double $\infty$-category.
% In \cite{Lou25}, Loubaton proves that double $\infty$-categories are equivalent to flagged $(\infty,2)$-categories.
% Using this identification, we unpack Haugseng's definition bellow

% \begin{definition}[\cite{Hau16}]
%     We have $\Cat{1}^{\V} \to \Mor_{1}^\beta(\V) \in \fCat{2}$ whose
%     \begin{enumerate}
%         \item \textbf{objects} are $\V$-categories,
%         \item \textbf{hom-categories} from $\A$ to $\B$ are given by the $\V$-enriched bifunctor category $[\A^\op \otimes \B, \V]_{\V}$.
%     \end{enumerate}
% \end{definition}

% In this section, we study its univalent completion and the notion of $\V$-absolute colimits.
% First, let's recall some facts about enriched ordinary category theory.

% \begin{proposition}
%     Suppose we have $\A, \B, \cat{C} \in \Cat{1}$, $U\in \A \to \B$, $L \in \begin{tikzcd}
%         \B & {\cat{C}}
%         \arrow[""{name=0, anchor=center, inner sep=0}, curve={height=-6pt}, from=1-1, to=1-2]
%         \arrow[""{name=1, anchor=center, inner sep=0}, curve={height=-6pt}, from=1-2, to=1-1]
%         \arrow["\dashv"{anchor=center, rotate=-90}, draw=none, from=0, to=1]
%     \end{tikzcd} \in R $, and $X$ a right adjoint to $LU$.
%     % https://q.uiver.app/#q=WzAsMyxbMCwwLCJcXEEiXSxbMiwwLCJcXEIiXSxbMSwxLCJcXEMiXSxbMSwwLCJVIiwyLHsiY3VydmUiOjF9XSxbMiwxLCJMIiwwLHsiY3VydmUiOi0xfV0sWzEsMiwiUiIsMCx7ImN1cnZlIjotMX1dLFswLDIsIlgiLDJdLFs0LDUsIiIsMSx7ImxldmVsIjoxLCJzdHlsZSI6eyJuYW1lIjoiYWRqdW5jdGlvbiJ9fV0sWzMsNiwiIiwwLHsibGV2ZWwiOjEsInN0eWxlIjp7Im5hbWUiOiJhZGp1bmN0aW9uIn19XV0=
% \[\begin{tikzcd}
% 	\A && \B \\
% 	& \cat{C}
% 	\arrow[""{name=0, anchor=center, inner sep=0}, "X"', from=1-1, to=2-2]
% 	\arrow[""{name=1, anchor=center, inner sep=0}, "U"', curve={height=6pt}, from=1-3, to=1-1]
% 	\arrow[""{name=2, anchor=center, inner sep=0}, "R", curve={height=-6pt}, from=1-3, to=2-2]
% 	\arrow[""{name=3, anchor=center, inner sep=0}, "L", curve={height=-6pt}, from=2-2, to=1-3]
% 	\arrow["\dashv"{anchor=center, rotate=-126}, draw=none, from=1, to=0]
% 	\arrow["\dashv"{anchor=center, rotate=-45}, draw=none, from=3, to=2]
% \end{tikzcd}\]
%     Then the following are equivalent we have a canonical $W_-$-weighted limit diagram of $X$ on the functor $B\Delta_+ \to \Cat{1}$ generated by the monad $RL$.
% \end{proposition}

% \begin{corollary}
%     Suppose that the adjunction $R \vdash L$ above was monadic.
%     Then we have a canonical lift $\bar{X}$ of $X$ through $R$.
%     Furthermore, one can write $\bar{X} = \Colim (LR)^{\circ k} \circ L \circ X$.
%     Because $U$ preserves colimit, we have $U\bar{X} = \Colim U(LR)^{\circ k} \circ L \circ X \to 1_{\A}$ canonical.
% \end{corollary}
% \begin{proof}
%     Cite Emily.
% \end{proof}

% \begin{proposition}
%     We have
%     % https://q.uiver.app/#q=WzAsMyxbMCwwLCJ7fV9BXFxNb2QoXFxWKSJdLFsyLDAsInt9X0JcXE1vZChcXFYpIl0sWzEsMSwiXFxWIl0sWzEsMCwiTSBcXG90aW1lc19CIC0iLDIseyJjdXJ2ZSI6MX1dLFsyLDEsIkIgXFxvdGltZXMgLSIsMCx7ImN1cnZlIjotMX1dLFsxLDIsIiIsMCx7ImN1cnZlIjotMX1dLFswLDIsIlt7fV9BTSwgLV0iLDJdLFs0LDUsIiIsMSx7ImxldmVsIjoxLCJzdHlsZSI6eyJuYW1lIjoiYWRqdW5jdGlvbiJ9fV0sWzMsNiwiIiwwLHsibGV2ZWwiOjEsInN0eWxlIjp7Im5hbWUiOiJhZGp1bmN0aW9uIn19XV0=
% \[\begin{tikzcd}
% 	{{}_A\Mod(\V)} && {{}_B\Mod(\V)} \\
% 	& \V
% 	\arrow[""{name=0, anchor=center, inner sep=0}, "{[M, -]}"', from=1-1, to=2-2]
% 	\arrow[""{name=1, anchor=center, inner sep=0}, "{M \otimes_B -}"', curve={height=6pt}, from=1-3, to=1-1]
% 	\arrow[""{name=2, anchor=center, inner sep=0}, curve={height=-6pt}, from=1-3, to=2-2]
% 	\arrow[""{name=3, anchor=center, inner sep=0}, "{B \otimes -}", curve={height=-6pt}, from=2-2, to=1-3]
% 	\arrow["\dashv"{anchor=center, rotate=-131}, draw=none, from=1, to=0]
% 	\arrow["\dashv"{anchor=center, rotate=-30}, draw=none, from=3, to=2]
% \end{tikzcd}\]
%     Then the following are equivalent we have a canonical $W_-$-weighted limit diagram of $X$ on the functor $B\Delta_+ \to \Cat{1}$ generated by the monad $RL$.
% \end{proposition}

% \subsection{Conventions on Higher Categories}

% \subsubsection{Many possible definitions}

% \autoref+
% There are many different appraoches to define higher categories, including:
% \begin{enumerate}
%     \item Rezk's complete Segal $\Theta_n$-spaces \cite{Rez10}, inspired by Joyal's unpublished ideas \cite{Joy97},
%         %functors out of "Joyal's cell or disk category" (parathensized pastings) to spaces,
%     \item Barwick's $n$-fold complete Segal spaces, inspired by Rezk's complete segal spaces (for $n=1$),
%         %functors out of grids to spaces,
%     \item Pellissier's categories Segal-enriched in any model category of $(\infty,n-1)$-categories,
%     \item Hirschowitz and Simpson's Segal $n$-categories inspired by the work of Tamsamani \cite{Tam96}, which are iterated versions of Pellisier's Segal-enrichment,
%         %functors out of grids (one more level) to sets,
%     \item Rezk-Kan's $n$-relative categories,
%     \item when $n=1$, Boardman-Vogt's quasicategories with the Joyal model structure,
%     \item when $n=1$, Lurie's marked simplicial sets,
%     \item when $n=1$, simplicial categories with Bergner's model structure,
%     \item when $n=1$, Ayala-Francis-Rozenblyum's striation sheaves,
%     \item when $n=2$, Lurie's scaled simplicial sets,
%     \item Gepner-Haugseng's $(\infty,1)$-categories enriched in any $(\infty,1)$-category of $(\infty,n-1)$-categories,
% \end{enumerate}

% There's also been many attempts at proving that all these definitions are equivalent:
% \begin{enumerate}
%     \item 
%     \item Haugseng \cite{Hau15} proves that the notion of enriched $(\infty,1)$-categories and Segal-enriched categories coïncide when the enrichment context is the $(\infty,1)$-category associated to a ``nice'' Cartesian model category.
%     This proves that 
% \end{enumerate}
% which all give equivalent $(\infty,1)$-categories by \cite{BS20}. Using this result, we will work with higher categories in a model-independent way.

% \subsubsection{Model-Independent Approach}

% We have a sequence of reflective and coreflective localizations of presentable $(\infty,1)$-categories

% % https://q.uiver.app/#q=WzAsNixbMSwwLCJcXENhdHsxfSJdLFsyLDAsIlxcQ2F0ezJ9Il0sWzMsMCwiXFxsZG90cyJdLFs0LDAsIlxcQ2F0e259Il0sWzUsMCwiXFxsZG90cyJdLFswLDAsIlxcU3BhIl0sWzEsMCwiIiwyLHsiY3VydmUiOjN9XSxbMSwwLCJ7LX1fe1xcbGVxIDF9IiwwLHsiY3VydmUiOi0zfV0sWzIsMSwiIiwyLHsiY3VydmUiOjN9XSxbMiwzLCIiLDAseyJzdHlsZSI6eyJ0YWlsIjp7Im5hbWUiOiJob29rIiwic2lkZSI6InRvcCJ9fX1dLFszLDQsIiIsMCx7InN0eWxlIjp7InRhaWwiOnsibmFtZSI6Imhvb2siLCJzaWRlIjoidG9wIn19fV0sWzIsMSwiey19X3tcXGxlcSAyfSIsMCx7ImN1cnZlIjotM31dLFszLDIsInstfV97XFxsZXEgbi0xfSIsMCx7ImN1cnZlIjotM31dLFs0LDMsInstfV97XFxsZXEgMH0iLDAseyJjdXJ2ZSI6LTN9XSxbNSwwLCIiLDAseyJzdHlsZSI6eyJ0YWlsIjp7Im5hbWUiOiJob29rIiwic2lkZSI6InRvcCJ9fX1dLFswLDUsInstfV97XFxsZXEgMH0iLDAseyJjdXJ2ZSI6LTN9XSxbMCw1LCIiLDAseyJjdXJ2ZSI6M31dLFswLDEsIiIsMix7InN0eWxlIjp7InRhaWwiOnsibmFtZSI6Imhvb2siLCJzaWRlIjoidG9wIn19fV0sWzEsMiwiIiwyLHsic3R5bGUiOnsidGFpbCI6eyJuYW1lIjoiaG9vayIsInNpZGUiOiJ0b3AifX19XSxbMywyLCIiLDAseyJjdXJ2ZSI6M31dLFs0LDMsIiIsMCx7ImN1cnZlIjozfV0sWzEwLDEzLCIiLDIseyJsZXZlbCI6MSwic3R5bGUiOnsibmFtZSI6ImFkanVuY3Rpb24ifX1dLFs5LDEyLCIiLDIseyJsZXZlbCI6MSwic3R5bGUiOnsibmFtZSI6ImFkanVuY3Rpb24ifX1dLFsxNCwxNSwiIiwwLHsibGV2ZWwiOjEsInN0eWxlIjp7Im5hbWUiOiJhZGp1bmN0aW9uIn19XSxbNiwxNywiIiwyLHsibGV2ZWwiOjEsInN0eWxlIjp7Im5hbWUiOiJhZGp1bmN0aW9uIn19XSxbMTYsMTQsIiIsMCx7ImxldmVsIjoxLCJzdHlsZSI6eyJuYW1lIjoiYWRqdW5jdGlvbiJ9fV0sWzgsMTgsIiIsMix7ImxldmVsIjoxLCJzdHlsZSI6eyJuYW1lIjoiYWRqdW5jdGlvbiJ9fV0sWzE3LDcsIiIsMix7ImxldmVsIjoxLCJzdHlsZSI6eyJuYW1lIjoiYWRqdW5jdGlvbiJ9fV0sWzE5LDksIiIsMix7ImxldmVsIjoxLCJzdHlsZSI6eyJuYW1lIjoiYWRqdW5jdGlvbiJ9fV0sWzE4LDExLCIiLDIseyJsZXZlbCI6MSwic3R5bGUiOnsibmFtZSI6ImFkanVuY3Rpb24ifX1dLFsyMCwxMCwiIiwyLHsibGV2ZWwiOjEsInN0eWxlIjp7Im5hbWUiOiJhZGp1bmN0aW9uIn19XV0=
% \[\begin{tikzcd}
% 	\Spa & {\Cat{1}} & {\Cat{2}} & \ldots & {\Cat{n}} & \ldots
% 	\arrow[""{name=0, anchor=center, inner sep=0}, hook, from=1-1, to=1-2]
% 	\arrow[""{name=1, anchor=center, inner sep=0}, "{{-}_{\leq 0}}", curve={height=-18pt}, from=1-2, to=1-1]
% 	\arrow[""{name=2, anchor=center, inner sep=0}, curve={height=18pt}, from=1-2, to=1-1]
% 	\arrow[""{name=3, anchor=center, inner sep=0}, hook, from=1-2, to=1-3]
% 	\arrow[""{name=4, anchor=center, inner sep=0}, curve={height=18pt}, from=1-3, to=1-2]
% 	\arrow[""{name=5, anchor=center, inner sep=0}, "{{-}_{\leq 1}}", curve={height=-18pt}, from=1-3, to=1-2]
% 	\arrow[""{name=6, anchor=center, inner sep=0}, hook, from=1-3, to=1-4]
% 	\arrow[""{name=7, anchor=center, inner sep=0}, curve={height=18pt}, from=1-4, to=1-3]
% 	\arrow[""{name=8, anchor=center, inner sep=0}, "{{-}_{\leq 2}}", curve={height=-18pt}, from=1-4, to=1-3]
% 	\arrow[""{name=9, anchor=center, inner sep=0}, hook, from=1-4, to=1-5]
% 	\arrow[""{name=10, anchor=center, inner sep=0}, "{{-}_{\leq n-1}}", curve={height=-18pt}, from=1-5, to=1-4]
% 	\arrow[""{name=11, anchor=center, inner sep=0}, curve={height=18pt}, from=1-5, to=1-4]
% 	\arrow[""{name=12, anchor=center, inner sep=0}, hook, from=1-5, to=1-6]
% 	\arrow[""{name=13, anchor=center, inner sep=0}, "{{-}_{\leq 0}}", curve={height=-18pt}, from=1-6, to=1-5]
% 	\arrow[""{name=14, anchor=center, inner sep=0}, curve={height=18pt}, from=1-6, to=1-5]
% 	\arrow["\dashv"{anchor=center, rotate=-87}, draw=none, from=0, to=1]
% 	\arrow["\dashv"{anchor=center, rotate=-93}, draw=none, from=2, to=0]
% 	\arrow["\dashv"{anchor=center, rotate=-90}, draw=none, from=3, to=5]
% 	\arrow["\dashv"{anchor=center, rotate=-90}, draw=none, from=4, to=3]
% 	\arrow["\dashv"{anchor=center, rotate=-97}, draw=none, from=6, to=8]
% 	\arrow["\dashv"{anchor=center, rotate=-83}, draw=none, from=9, to=10]
% 	\arrow["\dashv"{anchor=center, rotate=-83}, draw=none, from=7, to=6]
% 	\arrow["\dashv"{anchor=center, rotate=-97}, draw=none, from=12, to=13]
% 	\arrow["\dashv"{anchor=center, rotate=-97}, draw=none, from=11, to=9]
% 	\arrow["\dashv"{anchor=center, rotate=-83}, draw=none, from=14, to=12]
% \end{tikzcd}\]

% By abuse of notations, we don't distinguish an $(\infty,k)$-category $\cat{C} \in \Cat{k}$ from its image in $\Cat{n}$ for $n \geq k$.
% We let $\{\bullet\}$ be the terminal space, seen as the free category on a single object: it is the terminal object of $\Cat{n}$ for all $n$.
% We let $\emptyset$ denote the empty space: it's the initial object of $\Cat{n}$ for all $n$.
% By ``$c \in \cat{C}$'' we mean $c \in \{\bullet\} \to \cat{C}$.
% For $\cat{C} \in \Cat{n}$, the hom-categories can be seen as
% \[
%     \cat{C}_{\leq 0} \times \cat{C}_{\leq 0} \xto[\cat{C}(-,-)] \Cat{n-1}
% \]
% as is straightforward from the iterated-enrichment definition. We have a suspension functor $\sigma \in \Cat{n-1} \to \Cat{n}$ which sends an $(\infty, n-1)$-category $\cat{C}$ to the unique $(\infty, n)$-category with two objects $\{-,+\}$ and   
% \[\begin{cases}[rCl]
%     \sigma\cat{C}(+,+) &=& \{\bullet\}\\
%     \sigma\cat{C}(+,-) &=& \emptyset \\
%     \sigma\cat{C}(-,-) &=& \{\bullet\}\\
%     \sigma\cat{C}(-,+) &=& \cat{C}
% \end{cases}\]
% Equivalently, $\sigma\cat{C}$ is the object corepresenting the copresheaf
% \[\begin{cases}[rCl]
%     \Cat{n} &\longto& \Spa \\
%     \A &\longmapsto& \Lim \left(\A_{\leq 0} \times \A_{\leq 0} \xto[\A(-,-)] \Cat{n-1} \xto[{[\cat{C},-]}] \Spa \right)
% \end{cases}\]
% We let $c_0 = \{\bullet\}$ and for $n \geq 1$, let $c_n = \sigma^{n}c_0$ be the free category on a single $n$-cell.
% For all $n \geq 0$, we have an reflective localization
% % https://q.uiver.app/#q=WzAsMixbMSwwLCJcXENhdHtufSJdLFswLDAsIlxcdGV4dG5vcm1hbHtcXHRleHRiZntDYXR9fV97KG4sbil9Il0sWzEsMCwiIiwyLHsiY3VydmUiOjIsInN0eWxlIjp7InRhaWwiOnsibmFtZSI6Imhvb2siLCJzaWRlIjoidG9wIn19fV0sWzAsMSwiXFxob19uIiwyLHsiY3VydmUiOjJ9XSxbMywyLCIiLDIseyJsZXZlbCI6MSwic3R5bGUiOnsibmFtZSI6ImFkanVuY3Rpb24ifX1dXQ==
% \[\begin{tikzcd}
% 	{\textnormal{\textbf{Cat}}_{(n,n)}} & {\Cat{n}}
% 	\arrow[""{name=0, anchor=center, inner sep=0}, curve={height=12pt}, hook, from=1-1, to=1-2]
% 	\arrow[""{name=1, anchor=center, inner sep=0}, "{\ho_n}"', curve={height=12pt}, from=1-2, to=1-1]
% 	\arrow["\dashv"{anchor=center, rotate=-90}, draw=none, from=1, to=0]
% \end{tikzcd}\]
% where $\textnormal{\textbf{Cat}}_{(n,n)}$ is the full subcategory on the $n$-truncated objects.
% In particular, $\textnormal{\textbf{Cat}}_{(1,1)}$ is the usual $(2,1)$-category of ordinary categories and $\textnormal{\textbf{Cat}}_{(2,2)}$ is the usual $(3,1)$-category of bicategories.

% For $\cat{C}\in \Cat{n}$, we write $f\in a_{k-1} \to b_{k-1} \in \ldots \in a_0 \to b_0 \in \cat{C}$ to mean we have a diagram
% % https://q.uiver.app/#q=WzAsNixbNCwwLCJjX2siXSxbMywwLCJjX3trLTF9Il0sWzIsMCwiXFxsZG90cyJdLFsxLDAsImNfMSJdLFswLDAsIlxcYnVsbGV0Il0sWzYsMCwiXFxjYXR7Q30iXSxbNCwzLCJcXGRvbSIsMCx7ImN1cnZlIjotMX1dLFs0LDMsIlxcY29kIiwyLHsiY3VydmUiOjF9XSxbMywyLCJcXGRvbSIsMCx7ImN1cnZlIjotMX1dLFszLDIsIlxcY29kIiwyLHsiY3VydmUiOjF9XSxbMiwxLCJcXGRvbSIsMCx7ImN1cnZlIjotMX1dLFsyLDEsIlxcY29kIiwyLHsiY3VydmUiOjF9XSxbMSwwLCJcXGRvbSIsMCx7ImN1cnZlIjotMX1dLFsxLDAsIlxcY29kIiwyLHsiY3VydmUiOjF9XSxbMCw1LCJmIiwyLHsic3R5bGUiOnsiYm9keSI6eyJuYW1lIjoiZGFzaGVkIn19fV0sWzEsNSwiYl97ay0xfSIsMix7ImN1cnZlIjoyLCJzdHlsZSI6eyJib2R5Ijp7Im5hbWUiOiJkb3R0ZWQifX19XSxbMSw1LCJhX3trLTF9IiwwLHsiY3VydmUiOi0yLCJzdHlsZSI6eyJib2R5Ijp7Im5hbWUiOiJkb3R0ZWQifX19XSxbMyw1LCJiXzEiLDIseyJjdXJ2ZSI6Mywic3R5bGUiOnsiYm9keSI6eyJuYW1lIjoiZGFzaGVkIn19fV0sWzMsNSwiYV8xIiwwLHsiY3VydmUiOi0zLCJzdHlsZSI6eyJib2R5Ijp7Im5hbWUiOiJkb3R0ZWQifX19XSxbNCw1LCJiXzAiLDIseyJjdXJ2ZSI6NSwic3R5bGUiOnsiYm9keSI6eyJuYW1lIjoiZG90dGVkIn19fV0sWzQsNSwiYV8wIiwwLHsiY3VydmUiOi01LCJzdHlsZSI6eyJib2R5Ijp7Im5hbWUiOiJkb3R0ZWQifX19XV0=
% \[\begin{tikzcd}
% 	\bullet & {c_1} & \ldots & {c_{k-1}} & {c_k} && {\cat{C}}
% 	\arrow["\dom", curve={height=-6pt}, from=1-1, to=1-2]
% 	\arrow["\cod"', curve={height=6pt}, from=1-1, to=1-2]
% 	\arrow["{b_0}"', curve={height=30pt}, dotted, from=1-1, to=1-7]
% 	\arrow["{a_0}", curve={height=-30pt}, dotted, from=1-1, to=1-7]
% 	\arrow["\dom", curve={height=-6pt}, from=1-2, to=1-3]
% 	\arrow["\cod"', curve={height=6pt}, from=1-2, to=1-3]
% 	\arrow["{b_1}"', curve={height=18pt}, dashed, from=1-2, to=1-7]
% 	\arrow["{a_1}", curve={height=-18pt}, dotted, from=1-2, to=1-7]
% 	\arrow["\dom", curve={height=-6pt}, from=1-3, to=1-4]
% 	\arrow["\cod"', curve={height=6pt}, from=1-3, to=1-4]
% 	\arrow["\dom", curve={height=-6pt}, from=1-4, to=1-5]
% 	\arrow["\cod"', curve={height=6pt}, from=1-4, to=1-5]
% 	\arrow["{b_{k-1}}"', curve={height=12pt}, dotted, from=1-4, to=1-7]
% 	\arrow["{a_{k-1}}", curve={height=-12pt}, dotted, from=1-4, to=1-7]
% 	\arrow["f"', dashed, from=1-5, to=1-7]
% \end{tikzcd}\]


% \section{Higher Walking Adjunctions}

% \subsection{Preliminaries}

% \begin{definition}[Adjoint]
%     Let $n \geq 2$ and $\cat{C}\in \Cat{n}$.
%     We say that $l\in a \to b \in \cat{C}$ is a left adjoint if its equivalence class in $\ho_2(\cat{C})$ is a left adjoint.
%     For $k \geq 2$, say that $l\in a_{k-1} \to b_{k-1} \in \ldots \in a_0 \to b_0 \in \cat{C}$ is a left adjoint if its equivalence class in $\ho_2(\cat{C}(a_0,b_0)(a_1,b_1)\ldots(a_{k-2}, b_{k-2}))$ is a left adjoint.
% \end{definition}

% \begin{remark}[Being Adjoint is a Property]
%     In a bicategory $\B$, being an adjoint is a property of equivalence classes of elements in $[c_1,\B]$. 
%     As a consquence, in any $(\infty,n)$-category $\cat{C}$, the space of adjoint morphisms form a union of connected components of $[c_1, \cat{C}]$.
% \end{remark}

% Schanuel and Street \cite{SS86} defined combinatorially the free $(2,2)$-category $\Adj$ containing an adjunction $L \in \begin{tikzcd}
%         + & -
%         \arrow[""{name=0, anchor=center, inner sep=0}, curve={height=-6pt}, from=1-1, to=1-2]
%         \arrow[""{name=1, anchor=center, inner sep=0}, curve={height=-6pt}, from=1-2, to=1-1]
%         \arrow["\dashv"{anchor=center, rotate=-90}, draw=none, from=0, to=1]
%     \end{tikzcd} \in R $. We cite here two results fromn \cite{RV16} that we will use throughout this section.

% \begin{theorem}[Homotopy Uniqueness of Adjunctions {\cite[Thm. 4.4.11]{RV16}}]\label{them:riehl-verity}
%    Let $\cat{C}\in \Cat{2}$.
%    Precomposition by $L,R \in c_1 \to \Adj$ gives monomorphisms
%     % https://q.uiver.app/#q=WzAsMixbMCwwLCJbXFxBZGosXFxtYXRoY2Fse0N9XSJdLFsyLDAsIlxcbWF0aGNhbHtDfV5cXHRvIl0sWzAsMSwiTF4qIiwwLHsiY3VydmUiOi0xLCJzdHlsZSI6eyJ0YWlsIjp7Im5hbWUiOiJob29rIiwic2lkZSI6InRvcCJ9fX1dLFswLDEsIlJeKiIsMix7ImN1cnZlIjoxLCJzdHlsZSI6eyJ0YWlsIjp7Im5hbWUiOiJob29rIiwic2lkZSI6ImJvdHRvbSJ9fX1dXQ==
%     \[\begin{tikzcd}
%         {[\Adj,\mathcal{C}]} && {[c_1, \mathcal{C}]}
%         \arrow["{L^*}", curve={height=-6pt}, hook, from=1-1, to=1-3]
%         \arrow["{R^*}"', curve={height=6pt}, hook', from=1-1, to=1-3]
%     \end{tikzcd}\]
%     selecting the connected components of left (resp. right) adjoint morphisms.
% \end{theorem}

% \begin{corollary}
%     Let $n\geq 2$, $k\geq 1$ and $\cat{C}\in \Cat{n}$.
%     Precomposition by $\sigma^{k-1}L,\sigma^{k-1}R \in c_k \to \sigma^{k-1}\Adj$ gives monomorphisms
%      % https://q.uiver.app/#q=WzAsMixbMCwwLCJbXFxBZGosXFxtYXRoY2Fse0N9XSJdLFsyLDAsIlxcbWF0aGNhbHtDfV5cXHRvIl0sWzAsMSwiTF4qIiwwLHsiY3VydmUiOi0xLCJzdHlsZSI6eyJ0YWlsIjp7Im5hbWUiOiJob29rIiwic2lkZSI6InRvcCJ9fX1dLFswLDEsIlJeKiIsMix7ImN1cnZlIjoxLCJzdHlsZSI6eyJ0YWlsIjp7Im5hbWUiOiJob29rIiwic2lkZSI6ImJvdHRvbSJ9fX1dXQ==
%      \[\begin{tikzcd}
%          {[\sigma^{k-1}\Adj,\mathcal{C}]} && {[c_k, \mathcal{C}]}
%          \arrow["{\sigma^{k-1}L^*}", curve={height=-6pt}, hook, from=1-1, to=1-3]
%          \arrow["{\sigma^{k-1}R^*}"', curve={height=6pt}, hook', from=1-1, to=1-3]
%      \end{tikzcd}\]
%      selecting the connected components of left (resp. right) adjoint $k$-cells.
% \end{corollary}

% \begin{theorem}[Homotopy Coherent Lifting of Adjunctions {\cite[Thm. 4.3.11]{RV16}}]
%     Let $\cat{C}\in \Cat{2}$.
%     Suppose we have $(L, R, \eta, \epsilon)$ forming an adjunction in $\ho_2(\cat{C})$.
%     Then we have $\Adj \to \cat{C}$ lifting $(L, R, \eta, \epsilon)$.
% \end{theorem}

% \begin{corollary}
%     Let $\cat{C}\in \Cat{2}$.
%     Suppose we have $(L, R, \eta, \epsilon)$ forming an adjunction in ...
%     Then we have $\sigma^{k-1}\Adj \to \cat{C}$ lifting $(L, R, \eta, \epsilon)$.
% \end{corollary}\pablo{annoying to write}

% \begin{definition}
%     Let $n\geq 2$, $S \subset \{1,\ldots, n-1\}$ and $\cat{C} \in \Cat{n}$.
%     We say that $\cat{C}$ is \textbf{$S$-fully-adjoint} if for all $k \inset S$, the monomorphisms
%     % https://q.uiver.app/#q=WzAsMixbMCwwLCJbXFxBZGosXFxtYXRoY2Fse0N9XSJdLFsyLDAsIlxcbWF0aGNhbHtDfV5cXHRvIl0sWzAsMSwiTF4qIiwwLHsiY3VydmUiOi0xLCJzdHlsZSI6eyJ0YWlsIjp7Im5hbWUiOiJob29rIiwic2lkZSI6InRvcCJ9fX1dLFswLDEsIlJeKiIsMix7ImN1cnZlIjoxLCJzdHlsZSI6eyJ0YWlsIjp7Im5hbWUiOiJob29rIiwic2lkZSI6ImJvdHRvbSJ9fX1dXQ==
%     \[\begin{tikzcd}
%         {[\sigma^{k-1}\Adj,\mathcal{C}]} && {[c_k, \mathcal{C}]}
%         \arrow["{\sigma^{k-1}L^*}", curve={height=-6pt}, hook, from=1-1, to=1-3]
%         \arrow["{\sigma^{k-1}R^*}"', curve={height=6pt}, hook', from=1-1, to=1-3]
%     \end{tikzcd}\]
%     are equivalences.
%     We write $\Cat{n}^{S-\adj}$ for the full subcategory of $\Cat{n}$ on the $S$-fully-adjoint categories.
%     For $S = \{1,\ldots, n-1\}$, we write $\Cat{n}^{\adj} = \Cat{n}^{S-\adj}$.
% \end{definition}

% \begin{proposition}
%     Let $n\geq 2$ and $S \subset \{1,\ldots, n-1\}$. The category $\Cat{n}^{S-\adj}$ is locally finitely presentable and we have a (co)reflection
%     % https://q.uiver.app/#q=WzAsMixbMCwwLCJcXENhdHtufV57Uy1cXGFkan0iXSxbMiwwLCJcXENhdHtufSJdLFswLDEsIiIsMix7InN0eWxlIjp7InRhaWwiOnsibmFtZSI6Imhvb2siLCJzaWRlIjoidG9wIn19fV0sWzEsMCwiIiwyLHsiY3VydmUiOi00fV0sWzEsMCwiXFxhZGpuIiwyLHsiY3VydmUiOjR9XSxbMiwzLCIiLDIseyJsZXZlbCI6MSwic3R5bGUiOnsibmFtZSI6ImFkanVuY3Rpb24ifX1dLFs0LDIsIiIsMCx7ImxldmVsIjoxLCJzdHlsZSI6eyJuYW1lIjoiYWRqdW5jdGlvbiJ9fV1d
% \[\begin{tikzcd}
% 	{\Cat{n}^{S-\adj}} && {\Cat{n}}
% 	\arrow[""{name=0, anchor=center, inner sep=0}, hook, from=1-1, to=1-3]
% 	\arrow[""{name=1, anchor=center, inner sep=0}, curve={height=-24pt}, from=1-3, to=1-1]
% 	\arrow[""{name=2, anchor=center, inner sep=0}, "S-\adj"', curve={height=24pt}, from=1-3, to=1-1]
% 	\arrow["\dashv"{anchor=center, rotate=-101}, draw=none, from=0, to=1]
% 	\arrow["\dashv"{anchor=center, rotate=-79}, draw=none, from=2, to=0]
% \end{tikzcd}\]
%     For $S = \{1,\ldots, n-1\}$, we denote the reflection $n-\adj\in \Cat{n} \to \Cat{n}^{\adj}$.
% \end{proposition}
% \begin{proof}
%     Idea: use the complete segal space representation and explicitly build the adjoint. Then presentability follows.
% \end{proof}

% \begin{proposition}[Colimits have Adjoints]\label{prop:colim-adj}
%     Let $F \in \cat{I} \to \Cat{n}$ be a diagram.
%     We have that $\Colim F$ is $S$-fully-adjoint if and only if for all $i \in cat{I}$ and all $c_k \to Fi$ $k$-cell, the induced $k$-cell $c_k \to Fi \to \Colim F$ is $S$-fully-adjoint.
% \end{proposition}
% \begin{proof}
%     Idea: a $k$-cell in the colimit can always be written as a composition of $k$-cells in the diagram (and invertible data).
%     This is clear in the Grothendieck construction of the colimit, but hopefully also is clear via explicit definitions of higher categories (maybe some model category argument).
% \end{proof}

% \subsection{Higher Walking Adjunctions as Colimits}

% \subsubsection{Case $n = 2$}

% \begin{proposition}[Free Full Adjunction]
%     Denote $\mathbb{Z}^{\to\leftarrow}$ the poset whose objects are the integers and which has morphisms $2n \to 2n\pm 1$ for all $n \in \mathbb{Z}$.
%     We have a colimit diagram
%      % https://q.uiver.app/#q=WzAsMTAsWzQsMiwiXFxBZGoiXSxbMywzLCJjXzEiXSxbMiwyLCJcXEFkaiJdLFs2LDIsIlxcQWRqIl0sWzAsMiwiXFxjZG90cyJdLFs4LDIsIlxcY2RvdHMiXSxbMSwzLCJjXzEiXSxbNSwzLCJjXzEiXSxbNywzLCJjXzEiXSxbNCwwLCJcXEFkal5cXGluZnR5Il0sWzEsMCwiTCJdLFsxLDIsIlIiLDJdLFs3LDAsIlIiLDJdLFs2LDQsIlIiLDJdLFs2LDIsIkwiXSxbNywzLCJMIl0sWzgsMywiUiIsMl0sWzgsNSwiTCJdLFsyLDksIiIsMix7ImN1cnZlIjotMSwic3R5bGUiOnsiYm9keSI6eyJuYW1lIjoiZGFzaGVkIn19fV0sWzAsOSwiIiwyLHsic3R5bGUiOnsiYm9keSI6eyJuYW1lIjoiZGFzaGVkIn19fV0sWzMsOSwiIiwyLHsiY3VydmUiOjEsInN0eWxlIjp7ImJvZHkiOnsibmFtZSI6ImRhc2hlZCJ9fX1dLFs0LDksIiIsMix7ImN1cnZlIjotMiwic3R5bGUiOnsiYm9keSI6eyJuYW1lIjoiZGFzaGVkIn19fV0sWzUsOSwiIiwwLHsiY3VydmUiOjIsInN0eWxlIjp7ImJvZHkiOnsibmFtZSI6ImRhc2hlZCJ9fX1dXQ==
% \[\begin{tikzcd}
% 	&&&& {c_1^{2-\adj}} \\
% 	\\
% 	\cdots && \Adj && \Adj && \Adj && \cdots \\
% 	& {c_1} && {c_1} && {c_1} && {c_1}
% 	\arrow[curve={height=-12pt}, dashed, from=3-1, to=1-5]
% 	\arrow[curve={height=-6pt}, dashed, from=3-3, to=1-5]
% 	\arrow[dashed, from=3-5, to=1-5]
% 	\arrow[curve={height=6pt}, dashed, from=3-7, to=1-5]
% 	\arrow[curve={height=12pt}, dashed, from=3-9, to=1-5]
% 	\arrow["R"', from=4-2, to=3-1]
% 	\arrow["L", from=4-2, to=3-3]
% 	\arrow["R"', from=4-4, to=3-3]
% 	\arrow["L", from=4-4, to=3-5]
% 	\arrow["R"', from=4-6, to=3-5]
% 	\arrow["L", from=4-6, to=3-7]
% 	\arrow["R"', from=4-8, to=3-7]
% 	\arrow["L", from=4-8, to=3-9]
% \end{tikzcd}\]
% We denote this diagram as $C_1^{2-\adj}\in \mathbb{Z}^{\to\leftarrow} \to \Cat{2}$.
% \end{proposition}
% \begin{proof}
%    Firstly, we build the cocone diagram.
%    By definition, we know that $c_1^{2-\adj}$ is fully-adjoint and receives a universal functor $c_ 1 \to c_1^{2-\adj}$.
%    The space of cocones is the limit of the diagram $\mathbb{Z}^{\to\leftarrow} \to \Spa$ bellow.
%     % https://q.uiver.app/#q=WzAsNSxbMCwwLCJcXGxkb3RzIl0sWzIsMCwiW1xcQWRqLCBjXzFeezItXFxhZGp9XSJdLFs0LDAsIlxcY2RvdHMiXSxbMSwxLCJbY18xLCBjXzFeezItXFxhZGp9XSJdLFszLDEsIltjXzEsIGNfMV57Mi1cXGFkan1dIl0sWzAsMywiUl4qIl0sWzEsMywiTF4qIiwyXSxbMSw0LCJSXioiXSxbMiw0LCJMXioiLDJdXQ==
%     \[\begin{tikzcd}
%         \ldots && {[\Adj, c_1^{2-\adj}]} && \cdots \\
%         & {[c_1, c_1^{2-\adj}]} && {[c_1, c_1^{2-\adj}]}
%         \arrow["{R^*}", from=1-1, to=2-2]
%         \arrow["{L^*}"', from=1-3, to=2-2]
%         \arrow["{R^*}", from=1-3, to=2-4]
%         \arrow["{L^*}"', from=1-5, to=2-4]
%     \end{tikzcd}\]
%     By Theorem \ref{them:riehl-verity}, this is an infinite zig-zag of weak homotopy equivalences.
%     By picking the point in the zero-th space corresponding to the universal functor $c_ 1 \to c_1^{2-\adj}$, we get a canonical cocone diagram.

%     Secondly, we build a functor $c_1^{2-\adj} \to \Colim C_1^{2-\adj}$.
%     By universal property of $c_1^{2-\adj}$, it's sufficient to show that $\Colim C_1^{2-\adj} \inset \Cat{2}^{\adj}$ and to pick a $1$-cell $c_1 \to \Colim C_1^{2-\adj}$.
%     By Proposition \ref{prop:colim-adj}, it suffices to show that all $1$-cells in the diagram $C_1^{2-\adj}\in \mathbb{Z}^{\to\leftarrow} \to \Cat{2}$ are fully $1$-adjoint.
%     Note that each $1$-cell of each copy of $\Adj$ is a composition of the corresponding copies of $L$ and $R$, which in turn factors through one of the copies of $c_1$ in the diagram.
%     Hence, it is sufficient to show that each of the $c_1$ in the diagram are fully $1$-adjoint, which is clear by construction.
%     We now get a functor $c_1^{2-\adj} \to \Colim C_1^{2-\adj}$ by choosing the $1$-cell $c_1 \to \Colim C_1^{2-\adj}$ given by the universal cocone projection on the zero-th copy of $c_1$.
    
%     It's now clear that the two functors $c_1^{2-\adj} \to \Colim C_1^{2-\adj}$ and $\Colim C_1^{2-\adj} \to c_1^{2-\adj}$ are inverses to each other by examining the universal properties.
    
% \end{proof}

% Note that the diagram above is contractible.

% \subsubsection{Higher Cases}

% % \begin{definition}
% %     Let $n \geq 3$.
% %     Let $\mathcal{I} \subset \mathcal{P}(\{1,\ldots, n-1\})$.
% %     We say that $\mathcal{I}$ is \textbf{covering} if:
% %     \begin{enumerate}
% %         \item for all $I \inset \mathcal{I}$ and $J \subset I$, we have $J \inset \mathcal{I}$,
% %         \item for all $k \inset \{1,\ldots, n-1\}$, there exists $I \inset \mathcal{I}$ such that $k \inset I$.
% %     \end{enumerate}
% %     We denote by $\partial \mathcal{I}$ the set of maximal elements of $\mathcal{I}$.
% % \end{definition}

% % \begin{definition}
% %     Let $n \geq 3$ and $\mathcal{I} \subset \mathcal{P}(\{1,\ldots, n-1\})$ covering.
% %     We define a poset $\bar{\mathcal{I}}$ as follows:
% %     \begin{enumerate}
% %         \item
% %     \end{enumerate}
% % \end{definition}

% % \begin{theorem}
% %     Let $n \geq 3$ and $\mathcal{I} \subset \mathcal{P}(\{1,\ldots, n-1\})$ covering.
% %     We have $c_1^{n-\adj} = \Colim (\bar{\mathcal{I}} \to \Cat{n})$.
% % \end{theorem}
% % \begin{proof}
% %     We want to show that $\Colim (\bar{\mathcal{I}} \to \Cat{n})$ is fully-adjoint.
% %     By induction, the previous one is.
% %     We need to check the highest $\epsilon$ and $\eta$-cells we added are fully-adjoint.
% %     We check that they have two sided adjoints.
% % \end{proof}

% % \begin{definition}[Duality Datum]
% %     A $k$-duality data is a word of length $k$ using letters $\epsilon$ and $\eta$, with or without an overline bar.
% % \end{definition}

% % \begin{corollary}
% %     Let $n\geq 3$ and $\cat{C}$ be a higher category. The canonical
% %     \[\begin{tikzcd}
% %          {[c_1^{n-\adj} ,\mathcal{C}]} && {[c_1, \mathcal{C}]}
% %          \arrow[hook, from=1-1, to=1-3]
% %      \end{tikzcd}\]
% %     is a monomorphism. To show that some morphism $f$ is in the image connected components, it suffices to:
% %     \begin{enumerate}
% %         \item pick $\mathcal{I} \subset \mathcal{P}(\{1,\ldots, n-1\})$ covering,
% %         \item \ldots
% %     \end{enumerate}
% % \end{corollary}

% Fix $n \geq 3$.

% \begin{definition}[Poset of Duality Data]
%     We define a poset $\mathcal{P}_n$ of $2^{2n-2}-1$ elements as follows:
%     \begin{enumerate}
%         \item \textbf{Objects:} words of length at most $n$ with letters alternating between $\{\epsilon, \eta\}$ and $\{L, R\}$. If non-empty, the word is required to start with either $L$ or $R$.
%         \item \textbf{Inequalities:} we have $w \leq w'$ whenever $w$ ends $\epsilon$ or $\eta$ and $w'$ is obtained from $w$ by adding one letter.
%     \end{enumerate} 
% \end{definition}

% For example, for $n=3$, we get
% % https://q.uiver.app/#q=WzAsMTUsWzQsMywiXFxlbXB0eXNldCJdLFsyLDIsIkwiXSxbNiwyLCJSIl0sWzEsMSwiXFxlcHNpbG9uIEwiXSxbMywxLCJcXGV0YSBMIl0sWzUsMSwiXFxlcHNpbG9uIFIiXSxbNywxLCJcXGV0YSBSIl0sWzAsMCwiTCBcXGVwc2lsb24gTCJdLFsxLDAsIiBSXFxlcHNpbG9uIEwiXSxbMiwwLCJMIFxcZXRhIEwiXSxbMywwLCJSIFxcZXRhIEwiXSxbNSwwLCJMIFxcZXBzaWxvbiBSIl0sWzYsMCwiUiBcXGVwc2lsb24gUiJdLFs3LDAsIkwgXFxldGEgUiJdLFs4LDAsIlIgXFxldGEgUiJdLFswLDFdLFswLDJdLFs0LDFdLFszLDFdLFs1LDJdLFs2LDJdLFszLDddLFszLDhdLFs0LDldLFs0LDEwXSxbNSwxMV0sWzUsMTJdLFs2LDEzXSxbNiwxNF1d
% \[\begin{tikzcd}
% 	{L \epsilon L} & { R\epsilon L} & {L \eta L} & {R \eta L} && {L \epsilon R} & {R \epsilon R} & {L \eta R} & {R \eta R} \\
% 	& {\epsilon L} && {\eta L} && {\epsilon R} && {\eta R} \\
% 	&& L &&&& R \\
% 	&&&& \emptyset
% 	\arrow[from=2-2, to=1-1]
% 	\arrow[from=2-2, to=1-2]
% 	\arrow[from=2-2, to=3-3]
% 	\arrow[from=2-4, to=1-3]
% 	\arrow[from=2-4, to=1-4]
% 	\arrow[from=2-4, to=3-3]
% 	\arrow[from=2-6, to=1-6]
% 	\arrow[from=2-6, to=1-7]
% 	\arrow[from=2-6, to=3-7]
% 	\arrow[from=2-8, to=1-8]
% 	\arrow[from=2-8, to=1-9]
% 	\arrow[from=2-8, to=3-7]
% 	\arrow[from=4-5, to=3-3]
% 	\arrow[from=4-5, to=3-7]
% \end{tikzcd}\]

% \begin{definition}[Diagram $C_1^{n-\adj}$]
%     We define a diagram $C_1^{n-\adj} \in \mathcal{P}_n \to \Cat{n}$ as follows:
%     \begin{enumerate}
%         \item \textbf{Objects:} words ending in $\{L, R\}$ are sent to $\sigma^k \Adj$ where $2k+1$ is the length of the word; other words are sent to $c_k$ where again $2k+1$ is the length.
%         \item \textbf{Morphisms:} morphisms are labelled by the letter added to the word.
%     \end{enumerate}
% \end{definition}

% For example, for $n=3$, we get the following diagram in $\Cat{3}$:
% % https://q.uiver.app/#q=WzAsMTUsWzQsMywiY18xIl0sWzIsMiwiXFxBZGoiXSxbNiwyLCJcXEFkaiJdLFsxLDEsImNfMiJdLFszLDEsImNfMiJdLFs1LDEsImNfMiJdLFs3LDEsImNfMiJdLFswLDAsIlxcc2lnbWFcXEFkaiJdLFsxLDAsIlxcc2lnbWFcXEFkaiJdLFsyLDAsIlxcc2lnbWFcXEFkaiJdLFszLDAsIlxcc2lnbWFcXEFkaiJdLFs1LDAsIlxcc2lnbWFcXEFkaiJdLFs2LDAsIlxcc2lnbWFcXEFkaiJdLFs3LDAsIlxcc2lnbWFcXEFkaiJdLFs4LDAsIlxcc2lnbWFcXEFkaiJdLFswLDEsIkwiXSxbMCwyLCJSIiwyXSxbNCwxLCJcXGV0YSJdLFszLDEsIlxcZXBzaWxvbiIsMl0sWzUsMiwiXFxlcHNpbG9uIiwyXSxbNiwyLCJcXGV0YSJdLFszLDcsIkwiXSxbMyw4LCJSIiwyXSxbNCw5LCJMIl0sWzQsMTAsIlIiLDJdLFs1LDExLCJMIl0sWzUsMTIsIlIiLDJdLFs2LDEzLCJMIl0sWzYsMTQsIlIiLDJdXQ==
% \[\begin{tikzcd}
% 	{\sigma\Adj} & {\sigma\Adj} & {\sigma\Adj} & {\sigma\Adj} && {\sigma\Adj} & {\sigma\Adj} & {\sigma\Adj} & {\sigma\Adj} \\
% 	& {c_2} && {c_2} && {c_2} && {c_2} \\
% 	&& \Adj &&&& \Adj \\
% 	&&&& {c_1}
% 	\arrow["L", from=2-2, to=1-1]
% 	\arrow["R"', from=2-2, to=1-2]
% 	\arrow["\epsilon"', from=2-2, to=3-3]
% 	\arrow["L", from=2-4, to=1-3]
% 	\arrow["R"', from=2-4, to=1-4]
% 	\arrow["\eta", from=2-4, to=3-3]
% 	\arrow["L", from=2-6, to=1-6]
% 	\arrow["R"', from=2-6, to=1-7]
% 	\arrow["\epsilon"', from=2-6, to=3-7]
% 	\arrow["L", from=2-8, to=1-8]
% 	\arrow["R"', from=2-8, to=1-9]
% 	\arrow["\eta", from=2-8, to=3-7]
% 	\arrow["L", from=4-5, to=3-3]
% 	\arrow["R"', from=4-5, to=3-7]
% \end{tikzcd}\]

% \begin{proposition}
%     We have $c_1^{n-\adj} = \Colim C_1^{n-\adj}$
% \end{proposition}
% \begin{proof}
%     \bigpablo{First step: construct cocone (should be easy).
%     MAIN step: show the colimit is fully-adjunct (should come from colimits-have-adjoints and some non-trivial work on duality data).
%     Final step: inverse of one-another (should be easy too).
%     I'll do the second step first to get it out of the way.}
    
% \end{proof}

% \subsubsection{Redundancy of Duality Data}
% \bigpablo{Optional}

% \section{The Morita Category}

% % Finish facto in E_n section: should be easy to generalize Thm 2.25 (pushforward) from AFT17 to this setup
% % Don't lose time defining Mor_n: just use Haugseng's definition and result identifying the hom
% % Define the relevant stratified spaces (handle) and use pushforward to build them by tensoring
% % Finish the "center" section which identify some internal hom with some tensoring operation


% \subsection{Factorization Homology in $\mathcal{E}_n$-Monoïdal Categories}          

% \begin{proposition}
%     Let $\B \to \Bsc$ be a category of basics.
%     The functor $\Entr(-)=\Disk_{/-} \in \Mfld(\B) \to \Cat{1}$ is a homology theory.
% \end{proposition}
% \begin{proof}
%     By \cite[Thm. 1.1.(ii)]{GHN15}, we can identify the functor above with
%     % https://q.uiver.app/#q=WzAsMyxbMCwwLCJcXE1mbGQoXFxCKSJdLFsxLDAsIlxcUHNoKFxcRGlzayhcXEIpKSJdLFszLDAsIlxcQ2F0ezF9Il0sWzAsMSwiXFx5byJdLFsxLDIsIlxcbWF0aHNme2xheC19XFxDb2xpbSJdXQ==
% \[\begin{tikzcd}
% 	{\Mfld(\B)} & {\Psh(\Disk(\B))} && {\Cat{1}}
% 	\arrow["\yo", from=1-1, to=1-2]
% 	\arrow["{\mathsf{lax-}\Colim}", from=1-2, to=1-4]
% \end{tikzcd}\]
% From \cite[Cor. 2.28]{AFT17}, we know that the first functor - the restricted Yoneda functor of the inclusion $\Disk(\B) \mono \Mfld(\B)$ - is a homology theory, when $\Psh(\Disk(\B))$ is endowed with the Day convolution.
% Note that $\Disk_{/-} \in \Disk(\B) \to \Cat{1}$ is symmetric monoïdal.
% The second functor is obtained from it by the universal property of Day convolution, and hence preserves colimits and is symmetric monoïdal. 
% \end{proof}

% Recall from \cite[Constr. 2.4]{AFT17} that each stratified space $X$ defines an operad $\Disk_{/X}$.
% For $\V \in \EAlg{\infty}(\Cat)$, we denote $\Alg{X}(\V) = \Alg{{\Disk_{/X}}}(\V)$ the category of $\Disk_{/X}$-algebras in $\V$.
% Note that the operad $\Disk_{/\mathbb{R}^n}$ is the $\mathcal{E}_n$-operad.

% \begin{definition}
%     Let $\bar{X} \to X$ be a refinement of stratified spaces.
%     Let $\V \in \Alg{X}(\Cat{1})$.
%     We define an $\bar{X}$-algebra in $\V$ to be a symmetric monoïdal natural transformation
%     % https://q.uiver.app/#q=WzAsMyxbMCwwLCJcXERpc2tfey9cXGJhcntYfX0iXSxbMSwxLCJcXERpc2tfey9YfSJdLFsyLDAsIlxcQ2F0ezF9Il0sWzEsMiwiXFxWIiwyXSxbMCwyLCJcXEVudHIoLSkiLDAseyJjdXJ2ZSI6LTR9XSxbMCwxLCJcXHRleHRub3JtYWx7bG9jfSIsMl0sWzQsMSwiQSIsMCx7InNob3J0ZW4iOnsic291cmNlIjozMCwidGFyZ2V0IjozMH19XV0=
%     \[\begin{tikzcd}
%         {\Disk_{/\bar{X}}} && {\Cat{1}} \\
%         & {\Disk_{/X}}
%         \arrow[""{name=0, anchor=center, inner sep=0}, "{\Entr(-)}", curve={height=-24pt}, from=1-1, to=1-3]
%         \arrow["{\textnormal{loc}}"', from=1-1, to=2-2]
%         \arrow["\V"', from=2-2, to=1-3]
%         \arrow["A", between={0.3}{0.7}, Rightarrow, from=0, to=2-2]
%     \end{tikzcd}\]
% \end{definition}

% \bigpablo{Old version bellow}

% Recall that $\mathsf{D}^\fr_n \to \Bsc$ is defined to be the tangent classifyer $\tau_{\mathbb{R}^n}$.
% It can also be identified with the functor $\{\bullet\} \xto[\mathbb{R}^n] \Bsc$.

% \begin{definition}
%     We define $\Bim_n \to \Bsc$ to be the tangent classifyer $\tau_{\mathbb{R}^{n-1} \times C(S^0)}$.
% \end{definition}

% By \cite[Prop. 1.2.13]{AFT17a}, we have a localization functor $\Bim_n \to \mathsf{D}^\fr_n$.
% This induces a localization functor $\Disk(\Bim_n) \to \Disk(\mathsf{D}^\fr_n)$.
% This allows us to see any $\Disk(\mathsf{D}^\fr_n)$-algebra as a $\Disk(\Bim_n)$-algebra.

% One can also define $\Bim_n$ as follows.
% \bigpablo{Add definition of the category of framed $n$-manufolds with codim 1 normally framed submanifolds}
% \begin{definition}
%     Let \ldots
% \end{definition}

% \begin{definition}
%     Let $\V \in \EAlg{n}(\Cat{1})$.
%     We define a $\Disk(\Bim_n)$-algebra in $\V$ to be a symmetric monoïdal natural transformation
%     % https://q.uiver.app/#q=WzAsMyxbMCwwLCJcXERpc2soXFxCKSJdLFsxLDEsIlxcRGlzayhcXG1hdGhzZntEfV5cXGZyX24pIl0sWzIsMCwiXFxDYXR7MX0iXSxbMCwxXSxbMSwyLCJcXFYiXSxbMCwyLCJcXEVudHIoLSkiLDAseyJjdXJ2ZSI6LTR9XSxbNSwxLCJBIiwwLHsic2hvcnRlbiI6eyJzb3VyY2UiOjMwLCJ0YXJnZXQiOjMwfX1dXQ==
% \[\begin{tikzcd}
% 	{\Disk(\Bim_n)} && {\Cat{1}} \\
% 	& {\Disk(\mathsf{D}^\fr_n)}
% 	\arrow[""{name=0, anchor=center, inner sep=0}, "{\Entr(-)}", curve={height=-24pt}, from=1-1, to=1-3]
% 	\arrow[from=1-1, to=2-2]
% 	\arrow["\V", from=2-2, to=1-3]
% 	\arrow["A", shorten <=10pt, shorten >=10pt, Rightarrow, from=0, to=2-2]
% \end{tikzcd}\]
% \end{definition}

% Note that, by a similar reasonning to \cite[Prop. 4.8]{AFT17}, the data of such a $\Disk(\Bim_n)$-algebra is equivalent to the data of two $\mathcal{E}_n$-algebras $X$ and $Y$, and of a $\mathcal{E}_{n-1}$-algebra in $_{X}\Bim_Y$.

% By \cite[Thm. 2.43]{AFT17}, we can extend any such natural transformation to a morphism of homology theories
% \[
%     \Entr(-) \xto[A_*] \int_{-} \V
% \]

% Suppose additionally that we have $M^\mathsf{strat} \in \Mfld(\B)$.
% The functor of basics $\B \to \mathsf{D}^\fr_n$ \pablo{not true} gives $M$ a smoothing $M^\mathsf{man} \in \Mfld(\mathsf{D}^\fr_n)$ which is a smooth framed $n$-manifold.
% This smoothing also comes with a refinement map $M^\mathsf{strat} \to M^\mathsf{man}$ which induces an isomorphism
% \[
%     \int_{M^\mathsf{strat}} \V \simeq \int_{M^\mathsf{man}} \V 
% \]

% \begin{definition}
%     Using the notations above, suppose additionally that:
%     \begin{enumerate}
%         \item we have an embedding of framed manifold $\iota \in M^\mathsf{man} \mono \mathbb{R}^n$,
%         \item $\V$ has factors through $\Cat{1}^\sifted$, or in other words its $\mathcal{E}_n$-monoïdal tensor product $\otimes$ commutes with sifted colimits in each variables.
%     \end{enumerate}
%     We define the \textbf{factorization homology} of $M$ with coefficients in $A$ to be the colimit of the functor
%     \[
%         \Entr(M^\mathsf{strat}) \xto[A_*] \int_{M^\mathsf{strat}} \V \simeq \int_{M^\mathsf{man}} \V \xto[\iota_*] \int_{\mathbb{R}^n} \V \simeq \V
%     \]
%     where we identify $\V$ with the image of $\mathbb{R}^n$ under $\Disk(\mathsf{D}^\fr_n) \xto[\V] \Cat{1}$ by a slight abuse of notations.
% \end{definition}

% Clearly, we have a naturality in the choice of $\iota \in \Emb^\fr(M^\mathsf{man}, \mathbb{R}^n)$.

% This definition of factorization homology is functorial to some extent, but doesn't exhibit the full functoriality.
% Ideally, $M^\mathsf{strat}$ should be a point in some classifying space of stratifications of $\mathbb{R}^n$.
% The construction of such a space is given in \cite{AFR18} as ${}_{\mathbb{R}^n\backslash}\Mfld^{\sfr_n}$ in order to study the cobordism hypothesis.
% We will not need this construction here, as we only require some limited functoriality, which we spell out in the next few lemmas.

% \bigpablo{Add some excision lemma: I want $\int_{M^\mathsf{strat} \times \mathbb{R}^k} A$ to be $\mathcal{E}_k$}

% \begin{lemma}
%     We have a functor of basics $\Bim_n \to \mathsf{D}^\fr_n$.
% \end{lemma}

% \subsection{Inductive Definition}

% In \cite{Hau23}, Haugseng builds an inductive model for higher Morita categories, which we use here.
% He builds a functor
% \[
%     \Mor \in \EAlg{1}(\Cat{1}^\sifted) \to \fCat{2}
% \]
